\def\sode{3}
\begin{name}
	{\tenchude}
	{\tendethi}
	{\tentruong}
	{\thoigian}
\end{name}
\caulc
\Opensolutionfile{ans}[ans/ans-HXN-\sode-T]
%Câu hỏi
\begin{ex}%Câu 1
 Nguyên hàm của hàm số $f(x)=2^x$ là
 \choice
 {$\dfrac{2^{x+1}}{x+1}+C$}
 {\True $\dfrac{2^x}{\ln 2}+C$}
 {$\dfrac{2^x}{x}+C$}
 {$x{2^{x-1}}+C$}
 \loigiai{
 Chọn B.\\
 Ta có $\displaystyle\int\limits_{}^{}{2^x\text{d}x}=\dfrac{2^x}{\ln 2}+C$ .}
\end{ex}
\begin{ex}%Câu 2
 Cho hàm số $y=f(x)$ có bảng biến thiên như sau.
 \begin{center}
 \begin{tikzpicture}[>=stealth]
 \tkzTabInit[nocadre=false,lgt=1,espcl=2.5,deltacl=0.5]{$x$/.7 ,$y'$/.7,$y$/2}
 {$-\infty$ , $-2$ , $3$ , $+\infty$}
 \tkzTabLine{ , - , $0$ , + , $0$ , - , }
 \tkzTabVar{+/$+\infty$ , -/$-3$ , +/$2$ , -/$-\infty$}
 \end{tikzpicture}
 \end{center}
 Giá trị cực đại của hàm số đã cho bằng
 \choice
 {$3$}
 {\True $2$}
 {$-2$}
 {$-3$}
 \loigiai{
 Chọn B.\\
 Dựa vào bảng biến thiên, giá trị cực đại của hàm số là $y_{CT}=2$ .}
\end{ex}
\begin{ex}%Câu 3
 Cho cấp số cộng $\left(u_n\right)$ có $u_2=2;\,\,u_5=11$ . Công sai $d$ của cấp số cộng là
 \choice
 {1}
 {2}
 {4}
 {\True 3}
 \loigiai{
 Chọn D.\\
 Ta có: $\left\{\begin{aligned}
 &{u_2}=2\\ 
 &{u_5}=11\\ 
 \end{aligned}\right.\Leftrightarrow\left\{\begin{aligned}
 &{u_1}+d=2\\ 
 &{u_1}+4d=11\\ 
 \end{aligned}\right.\Leftrightarrow\left\{\begin{aligned}
 &{u_1}=-1\\ 
 & d=3\\ 
 \end{aligned}\right.$ . Vậy công sai của cấp số cộng là $d=3$ .}
\end{ex}
\begin{ex}%Câu 4
 Cho hàm số $y=f(x)$ có đạo hàm $f'(x)=x+1$ với mọi $x\in\mathbb{R}$ . Hàm số đã cho nghịch biến trên khoảng nào dưới đây?
 \choice
 {$\left(-1\,;\,\,+\infty\right)$}
 {$\left(1\,;\,\,+\infty\right)$}
 {\True $\left(-\infty\,;\,\,-1\right)$}
 {$\left(-\infty\,;\,\,1\right)$}
 \loigiai{
 Chọn C.\\
 }
\end{ex}
\begin{ex}%Câu 5
 Trong không gian $Oxyz$ , đường thẳng $d$ đi qua điểm $M\left(1\,;\,\,-1\,;\,\,3\right)$ và song song với đường thẳng$\Delta :\dfrac{x-2}{2}=\dfrac{y+1}{1}=\dfrac{z+3}{-1}$ có phương trình là
 \choice
 {\True $\left\{\begin{aligned}
 & x=1+2t\\ 
 & y=-1+t\\ 
 & z=3-t\\ 
 \end{aligned}\right.$}
 {$\left\{\begin{aligned}
 & x=1+2t\\ 
 & y=-1+t\\ 
 & z=3+t\\ 
 \end{aligned}\right.$}
 {$\left\{\begin{aligned}
 & x=2+t\\ 
 & y=1-t\\ 
 & z=-1+3t\\ 
 \end{aligned}\right.$}
 {$\left\{\begin{aligned}
 & x=1+2t\\ 
 & y=1+t\\ 
 & z=3-t\\ 
 \end{aligned}\right.$}
 \loigiai{
 Chọn A.\\
 Đường thẳng d song song $\Delta :\dfrac{x-2}{2}=\dfrac{y+1}{1}=\dfrac{z+3}{-1}$ nên có vectơ chỉ phương $\vec{u}=\left(2\,;\,\,1\,;\,\,-1\right)$ ; mà d qua $M\left(1\,;\,\,-1\,;\,\,3\right)$ nên có phương trình tham số $\left\{\begin{aligned}
 & x=1+2t\\ 
 & y=-1+t\\ 
 & z=3-t\\ 
 \end{aligned}\right.$ .}
\end{ex}
\begin{ex}%Câu 6
 Tập nghiệm của bất phương trình $\log_{\dfrac{1}{2}}\left(9-x^2\right)<0$ chứa bao nhiêu số nguyên?
 \choice
 {1}
 {\True $5$}
 {$3$}
 {4}
 \loigiai{
 Chọn B.\\
 Ta có: $\log_{\dfrac{1}{2}}\left(9-x^2\right)<0\Leftrightarrow 9-x^2>\left(\dfrac{1}{2}\right)^0\Leftrightarrow 9-x^2>1\Leftrightarrow{x^2}<8\Leftrightarrow-2\sqrt{2}<x<2\sqrt{2}$ .\\
 Tập nghiệm bất phương trình chứa 5 số nguyên là: $-2\,;\,\,1\,;\,\,0\,;\,\,1\,;\,\,2$ .}
\end{ex}
\begin{ex}%Câu 7
 Trong không gian $Oxyz$ , cho hai vectơ $\vec{u}=\left(1\,;\,\,-4\,;\,\,0\right)$ và $\vec{v}=\left(-1\,;\,\,-2\,;\,\,1\right)$ . Vectơ $\vec{u}+3\vec{v}$ có tọa độ là
 \choice
 {\True $\left(-2\,;\,\,-10\,;\,\,3\right)$}
 {$\left(-2\,;\,\,-6\,;\,\,3\right)$}
 {$\left(-4\,;\,\,-8\,;\,\,4\right)$}
 {$\left(-2\,;\,\,-10\,;\,\,-3\right)$}
 \loigiai{
 Chọn A.\\
 Ta có: $\vec{u}+3\vec{v}=\left(1\,;\,\,-4\,;\,\,0\right)+3\left(-1\,;\,\,-2\,;\,\,1\right)=\left(-2\,;\,\,-10\,;\,\,3\right)$ .}
\end{ex}
\begin{ex}%Câu 8
 Tâm đối xứng của đồ thị hàm số $y=\dfrac{3x+1}{x-2}$ có tọa độ là
 \choice
 {$\left(3\,;\,\,-2\right)$}
 {$\left(3\,;\,\,2\right)$}
 {$\left(-2\,;\,\,3\right)$}
 {\True $\left(2\,;\,\,3\right)$}
 \loigiai{
 Chọn D.\\
 Đồ thị hàm số có tiệm cận đứng $x\,=2$ và tiệm cận ngang $y\,=\,3$ .\\
 Do đó tâm đối xứng của đồ thị hàm số là $\left(2\,;\,\,3\right)$ .}
\end{ex}
\begin{ex}%Câu 9
 Cho mẫu số liệu ghép nhóm ở bảng sau. Khoảng tứ phân vị của mẫu số liệu ghép nhóm (làm tròn đến hàng phần trăm) là
 \choice
 {$19,15$}
 {$21,32$}
 {$20,07$}
 {\True $22,23$}
 \loigiai{
 Chọn D.\\
 Tứ phân vị thứ nhất của mẫu số liệu gốc là $\dfrac{x_6+x_7}{2}\in\left[30\,;\,\,40\right)$ nên tứ phân vị thứ nhất của mẫu số liệu ghép nhóm là $Q_1=30+\dfrac{\dfrac{25}{4}-3}{7}\cdot 10=\dfrac{485}{14}$ .\\
 Tứ phân vị thứ ba của mẫu số liệu gốc là $\dfrac{x_{19}+x_{20}}{2}\in\left[50\,;\,\,60\right)$ nên tứ phân vị thứ ba của mẫu số liệu ghép nhóm là $Q_3=50+\dfrac{3\cdot\dfrac{25}{4}-16}{4}\cdot 10=\dfrac{455}{8}$ .\\
 Khoảng tứ phân vị của mẫu số liệu ghép nhóm trên là $Q_3\,-\,Q_1\,=\,\dfrac{1245}{56}\approx\,22,23$ .}
\end{ex}
\begin{ex}%Câu 10
 Cho hàm số $f(x)=x^2+\sin x+1$ . Biết rằng $F(x)$ là một nguyên hàm của hàm số $f(x)$ và thỏa mãn $F(0)=1$ . Khi đó $F(x)$ bằng
 \choice
 {$F(x)=x^3-\cos x+x+2$}
 {\True $F(x)=\dfrac{x^3}{3}-\cos x+x+2$}
 {$F(x)=\dfrac{x^3}{3}+\cos x+x$}
 {$F(x)=\dfrac{x^3}{3}+\cos x+2$}
 \loigiai{
 Chọn B.\\
 Ta có $F(x)=\displaystyle\int{\left(x^2+\sin x+1\right)}\text{d}x=\dfrac{x^3}{3}-\cos x+x+C$ .\\
 Theo giả thiết: $F(0)=1\Rightarrow\dfrac{0^3}{3}-\cos 0+0+C=1\Rightarrow C=2$ .\\
 Do đó $F(x)=\dfrac{x^3}{3}-\cos x+x+2$ .}
\end{ex}
\begin{ex}%Câu 11
 Người ta thống kê lại đường kính thân gỗ của một số cây xoan đào 6 năm tuổi được trồng ở một lâm trường ở bảng sau:\\
 \centerline{\begin{tabular}{|c|c|c|c|c|c|}
 \hline
 Đường kính $(cm)$ & $[40\,;\,\,45)$ & $[45\,;\,\,50)$ & $[50\,;\,\,55)$ & $[55\,;\,\,60)$ & $[60\,;\,\,65)$\\
 \hline
 Tần số & 5 & 20 & 18 & 7 & 3\\
 \hline
 \end{tabular}}\\
 Khoảng biến thiên của mẫu số liệu ghép nhóm trên là
 \choice
 {\True 25}
 {30}
 {6}
 {69,8}
 \loigiai{
 Chọn A.\\
 Khoảng biến thiên của mẫu số liệu ghép nhóm trên là $65-40=25\,\,\,cm.$}
\end{ex}
\begin{ex}%Câu 12
 Cho hình hộp $ABCD.A'{B}'{C}'{D}'$ (xem hình vẽ). Phát biểu nào sau đây là đúng? 
 \choice
 {$\overrightarrow{A{C}'}=\overrightarrow{AB}+\overrightarrow{A{B}'}+\overrightarrow{AD}$}
 {\True $\overrightarrow{D{B}'}=\overrightarrow{DA}+\overrightarrow{D{D}'}+\overrightarrow{DC}$}
 {$\overrightarrow{A{C}'}=\overrightarrow{AC}+\overrightarrow{AB}+\overrightarrow{AD}$}
 {$\overrightarrow{DB}=\overrightarrow{DA}+\overrightarrow{D{D}'}+\overrightarrow{DC}$}
 \loigiai{
 Chọn B.\\
 Theo quy tắc hình hộp ta có$\overrightarrow{D{B}'}=\overrightarrow{DA}+\overrightarrow{D{D}'}+\overrightarrow{DC}$ .
 }
 \end{ex}
\Closesolutionfile{ans}
\cauds
\Opensolutionfile{ans}[ans/ans-HXN-\sode-TF]
%Câu hỏi
\begin{ex}
    \immini[thm]{Xét tam giác $ABC$ có $AC = 2AB$ và $BC = 10\ cm$. Trên cạnh $AC$ lấy điểm $D$ sao cho $AD = \dfrac{1}{4}AC$, trên cạnh $AB$ lấy điểm $E$ sao cho $AE = \dfrac{1}{4}AB$, trên cạnh $AD$ lấy điểm $F$ sao cho $AF = \dfrac{1}{4}AD$ và tiếp tục lấy các điểm $G, H, I, J, \dots$ (vô hạn lần) theo quy luật đó.
        
        \choiceTF
        {\True $\dfrac{AB}{AC} = \dfrac{AD}{AB}$}
        {Tam giác $ABD$ đồng dạng với tam giác $ABC$}
        {$BD = 5\ cm$; $DE = 3\ cm$}
        {\True Độ dài đường gấp khúc $CBDEFGH\dots$ bằng $20\ cm$}}{\includegraphics[scale=1.5]{img/HXN-3.13}}
    \loigiai{
    \begin{itemchoice}
        \itemch Ta có: $AD = \dfrac{1}{4}AC = \dfrac{1}{4} \cdot 2AB = \dfrac{1}{2}AB$; suy ra $\dfrac{AD}{AB} = \dfrac{1}{2}$. $\dfrac{AB}{AC} = \dfrac{AB}{2AB} = \dfrac{1}{2}$. Suy ra $\dfrac{AD}{AB} = \dfrac{AB}{AC}$.
        \itemch Hai tam giác $ABD$ và $ACB$ đồng dạng vì có góc $\hat{A}$ chung và $\dfrac{AB}{AC} = \dfrac{AD}{AB}$.
        \itemch Từ câu b) ta suy ra $\dfrac{AB}{AC} = \dfrac{AD}{AB} = \dfrac{BD}{CB} = \dfrac{1}{2} \Rightarrow BD = \dfrac{1}{2}BC = \dfrac{1}{2} \cdot 10 = 5\ cm$. Hoàn toàn tương tự, ta chứng minh được hai tam giác $ADB$ và $AED$ đồng dạng, suy ra $DE = \dfrac{1}{2}BD = \dfrac{1}{2} \cdot 5 = 2,5\ cm$.
        \itemch Độ dài đường gấp khúc $CBDEFGH\dots$ bằng $l = CB + BD + DE + EF + FG + \dots = 10 + 5 + 2,5 + \dots$ Đây là tổng của một cấp số nhân lùi vô hạn có số hạng đầu $u_1 = 10$, công bội $q = \dfrac{1}{2}$. Do đó $l = \dfrac{u_1}{1-q} = \dfrac{10}{1-\dfrac{1}{2}} = \dfrac{10}{\dfrac{1}{2}} = 20\ cm$.
        \end{itemchoice}}
\end{ex}

\begin{ex}
    \immini
    {
        Xét một hệ trục tọa độ $Oxyz$ được cho sẵn, đơn vị trên mỗi trục là dm, mặt ngoài của một quả bóng được mô hình hóa bởi phương trình mặt cầu $(x-2)^2+(y+1)^2+(z+1)^2=6$, quả bóng nằm yên trên sàn nhà. Người ta nhìn thấy một tấm ván ngã xuống đè lên quả bóng, phần giao của tấm ván và sàn nhà là đường thẳng $d$ có phương trình $\dfrac{x+2}{2}=\dfrac{y+1}{-3}=\dfrac{z}{1}$.  Gọi $A$, $B$ lần lượt là hai tiếp điểm của tấm ván, sàn nhà với quả bóng và $I$ là tâm quả bóng.
        \choiceTF
    {\True Quả bóng có tâm $I(2;-1;-1)$ và bán kính $R=\sqrt{6}$}
    {Khoảng cách từ tâm quả bóng đến đường thẳng d bằng $2\sqrt{6}$}
    {\True Nếu $\cos\widehat{AIB}$ bằng $\dfrac{a}{b}$ (phân số tối giản) thì giá trị $a^2+b^2=82$}
    {\True Một con kiến bò từ vị trí $A$ đến vị trí $B$ trên quả bóng với tốc độ $2$ cm/s; thời gian ngắn nhất cho chuyến đi này là $21$ giây (làm tròn đến hàng đơn vị)}
    }
    {
        \includegraphics[width=5cm]{img/HXN-3-14}
    }
    \loigiai{
        \begin{itemchoice}
            \itemch Mặt ngoài quả bóng là mặt cầu $(S)$ có tâm $I(2;-1;-1)$ và bán kính $R=\sqrt{6}$.
            \itemch Đường thẳng  $d\colon \dfrac{x+2}{2}=\dfrac{y+1}{-3}=\dfrac{z}{1}$ qua $A(-2;-1;0)$ và có VTCP $\vec{u}_d=(2;-3;1)$.\\
            Ta có $\vec{AI}=(4;0;-1);\left[{\vec{u}}_d,\vec{AI}\right]=(3;6;12)$.\\
            Do đó $d\left(I,d\right)=\dfrac{\left| \left[{\vec{u}}_d,\vec{AI}\right] \right|}{\left| {\vec{u}}_d \right|}=\dfrac{\sqrt{3^2+6^2+12^2}}{\sqrt{2^2+(-3)^2+1^2}}=\dfrac{3\sqrt{6}}{2}$.
            \immini
            {
            \itemch Gọi $K$ là hình chiếu của $I$ trên $d$ thì $KI=\dfrac{3\sqrt{6}}{2}$ và $KA\perp IA$;\\
            suy ra $\cos\widehat{AIK}=\dfrac{IA}{IK}=\dfrac{2}{3}$\\
            Do vậy $\cos\widehat{AIB}=2\cos^2\widehat{AIK}-1 =-\dfrac{1}{9}=\dfrac{a}{b}$\\
            $\Rightarrow a^2+b^2=82$.
            \itemch Độ dài cung tròn bé nhất mà con kiến có thể đi: $$l_{\wideparen{AB}}=R\times \widehat{AIB}=\sqrt{6}\times \arccos \left(-\dfrac{1}{9}\right)\approx 4{,}12\,dm$$
            Thời gian tối thiểu để kiến đến nơi là $\dfrac{l_{\wideparen{AB}}\times 10}{2}\approx 21$ giây.
            }
            {
                \includegraphics[width=5cm]{img/HXN-3-14-LG}
            }
        \end{itemchoice}
    }
\end{ex}
\begin{ex}
Thám tử lừng danh Sherlock Holmes đang điều tra một vụ án được thực hiện độc lập bởi một trong hai nghi phạm là McFarlane và Oldacre. Ban đầu thám tử đã có bằng chứng ngang nhau chống lại cả hai người.\\
Trong quá trình điều tra thêm tại hiện trường vụ án, Sherlock Holmes phát hiện rằng thủ phạm có nhóm máu mà chỉ $10\%$ dân số có; và Oldacre có nhóm máu này, còn nhóm máu của McFarlane thì chưa biết.\\
Gọi $A$ là biến cố: \lq\lq McFarlane là thủ phạm\rq\rq; $B$ là biến cố: \rq\rq Oldacre là thủ phạm\rq\rq; $C$ là biến cố: \lq\lq Nhóm máu của nghi phạm trùng với nhóm máu thủ phạm thực sự\rq\rq.
\choiceTF
{Trong quá trình điều tra, nếu Sherlock Holmes biết chắc chắn McFarlane không là  thủ phạm, khi đó xác suất để Oldacre là thủ phạm bằng $0{,}98$}
{$P\left(C|A\right)=0{,}1;P\left(C|B\right)=0{,}5$}
{Dựa trên thông tin về nhóm máu, xác suất để Oldacre là thủ phạm bằng $\dfrac{9}{11}$}
{\True Dựa trên thông tin về nhóm máu, xác suất để McFarlane cũng có nhóm máu trùng với nhóm máu thủ phạm bằng $\dfrac{2}{11}$}
\loigiai{
    \begin{itemchoice}
        \itemch Ta có: $P\left(B|\bar{A}\right)=\dfrac{P\left(B\bar{A}\right)}{P\left({\bar{A}}\right)}=\dfrac{P(B)}{P\left({\bar{A}}\right)}=\dfrac{\dfrac{1}{2}}{\dfrac{1}{2}}=1$. \\
        Như vậy nếu McFarlane không là  thủ phạm thì chắc chắn Oldacre là thủ phạm.
        \itemch Ta có: $P\left(C|A\right)=0{,}1;P\left(C|B\right)=1$.
        \itemch Áp dụng công thức Bayes:$P\left(B\mid C\right)=\dfrac{P(B)\cdot P\left(C\mid B\right)}{P(B)\cdot P\left(C\mid B\right)+P(A)\cdot P\left(C\mid A\right)}=\dfrac{\dfrac{1}{2}\cdot 1}{\dfrac{1}{2}\cdot 1+\dfrac{1}{2}\cdot \dfrac{1}{10}}=\dfrac{10}{11}$.
        \itemch Từ câu c) ta có $P\left(\bar{B}|C\right)=1-P\left(B\mid C\right)=\dfrac{1}{11}$.
        Gọi $D$ là biến cố \lq\lq McFarlane có cùng nhóm máu với thủ phạm biết rằng Oldacre có cùng nhóm máu với thủ phạ\rq\rq.
        Ta có $P(D)=\dfrac{10}{11}\cdot \dfrac{1}{10}+\dfrac{1}{11}\cdot 1=\dfrac{2}{11}$.
        (Nếu Oldacre là thủ phạm, xác suất để McFarlane có cùng nhóm máu với thủ phạm bằng $\dfrac{1}{10}$; nếu Oldacre không là thủ phạm thì xác suất để McFarlane có cùng nhóm máu với thủ phạm bằng $1$).
    \end{itemchoice}
}
\end{ex}
\begin{ex}
\immini
{
    Cho $y=f(x)$, $y=g(x)$ lần lượt là các hàm đa thức bậc ba và bậc nhất có đồ thị như hình vẽ.
    Biết diện tích hình $S$ (được tô màu) bằng $\dfrac{250}{81}$. 
    \choiceTF
    {Hàm số $g(x)=\dfrac{3}{5}x-\dfrac{1}{5}$}
    {\True Hàm số $\displaystyle\int\limits_{-2}^{\tfrac{4}{3}}{\left[f(x)-g(x)\right]\mathrm{\,d}x}=\dfrac{250}{81}$}
    {Hàm số $f(x)=\dfrac{3}{10}(x+2)\left(x-\dfrac{4}{3}\right)(x-3)+\dfrac{3}{5}x+\dfrac{1}{5}$}
    {$\int\limits_0^2{f(x)\mathrm{\,d}x}=\dfrac{37}{15}$}
}
{
    \includegraphics[width=5cm]{img/HXN-3-16}
}
    \loigiai{
        \begin{itemchoice}
            \itemch Ta có $g(x)$ là hàm số bậc nhất đi qua các điểm $A\left(\dfrac{4}{3};1\right)$, $B(3;2)$ nên $g(x)=\dfrac{3}{5}x+\dfrac{1}{5}$.
            \itemch Ta thấy hai đồ thị hàm số $y=f(x),y=g(x)$ cắt nhau tại điểm có $y=-1$; thay vào đường thẳng $y=\dfrac{3}{5}x+\dfrac{1}{5}$  thì $x=-2$.
            Do đó $S=\displaystyle\int\limits_{-2}^{\tfrac{4}{3}}{\left[f(x)-g(x)\right]\mathrm{\,d}x}=\dfrac{250}{81}$.
            \itemch Đặt $f(x)-g(x)=a(x+2)\left(x-\dfrac{4}{3}\right)(x-3)$ với $a>0$.\\
            Ta có: $S=\int\limits_{-2}^{\tfrac{4}{3}}{\left[f(x)-g(x)\right]\mathrm{\,d}x}\Leftrightarrow \displaystyle\int\limits_{-2}^{\tfrac{4}{3}}{\left[a(x+2)\left(x-\dfrac{4}{3}\right)(x-3)\right]\mathrm{\,d}x}=\dfrac{250}{81}\Leftrightarrow a=\dfrac{3}{20}$.\\
            Khi đó $f(x)-g(x)=\dfrac{3}{20}(x+2)\left(x-\dfrac{4}{3}\right)(x-3)\Leftrightarrow f(x)=\dfrac{3}{20}(x+2)\left(x-\dfrac{4}{3}\right)(x-3)+\dfrac{3}{5}x+\dfrac{1}{5}$.
            \itemch $\displaystyle\int\limits_0^2{f(x)\mathrm{\,d}x}=\displaystyle\int\limits_0^2{\left[\dfrac{3}{20}(x+2)\left(x-\dfrac{4}{3}\right)(x-3)+\dfrac{3}{5}x+\dfrac{1}{5}\right]\mathrm{\,d}x}=\dfrac{34}{15}$.
        \end{itemchoice}
    }
\end{ex}
\Closesolutionfile{ans}
\caukq
\Opensolutionfile{ans}[ans/ans-HXN-\sode-SA]
%Câu hỏi
\begin{ex}%Câu 17
 Cho tứ diện đều ABCD có tất cả cạnh bằng 2. Tính khoảng cách của hai đường thẳng chéo nhau AB và CD (làm tròn đến hàng phần trăm).
 \shortans{1,41}
\loigiai{
 \begin{center}
 \includegraphics[scale=1]{img/HXN-3.17a}
 \end{center}
 Gọi I, J theo thứ tự là trung điểm của AB, CD.\\
 Các tam giác ABC, ABD đều có I là trung điểm AB nên\\ $\left\{\begin{aligned}
 & AB\perp CI\\ 
 & AB\perp DI\\ 
 \end{aligned}\right.\Rightarrow AB\perp\left(ICD\right)$ mà $IJ\subset\left(ICD\right)\Rightarrow AB\perp IJ$ (1).\\
 Tương tự, các tam giác ACD, BCD đều có J là trung điểm CD nên\\ $\left\{\begin{aligned}
 & CD\perp AJ\\ 
 & CD\perp BJ\\ 
 \end{aligned}\right.\Rightarrow CD\perp\left(ABJ\right)$ , mà $IJ\subset\left(JAB\right)\Rightarrow CD\perp IJ$ (2).\\
 Từ (1) và (2) suy ra IJ là đoạn vuông góc chung của hai đường thẳng AB, CD.\\
 Ta có: $CI=\dfrac{2\sqrt{3}}{2}=\sqrt{3}$ ; $IJ=\sqrt{C{I^2}-C{J^2}}=\sqrt{3-1}=\sqrt{2}\approx 1,41$ .\\
 Vậy khoảng cách hai đường thẳng AB, CD xấp xỉ 1,41.}
 \end{ex}
 \begin{ex}%Câu 18
Buổi họp mặt cuối năm của VFF diễn ra trong không khí hân hoan phấn khởi sau khi ĐTQG Việt Nam vô địch AFF cup 2024. Vào cuối buổi họp thì HLV Kim Sang-sik chỉ bắt tay với một số người, còn lại tất cả thành viên đều bắt tay với nhau, hai người bất kì thì bắt tay không quá một lần. Hỏi cuộc họp này có bao nhiêu người tham dự biết rằng đã có tổng cộng 2014 cái bắt tay được thực hiện?
\shortans{64}
\loigiai{
Gọi n là người có mặt trong cuộc họp $\left(n\in\mathbb{N}\right)$ .\\
Số cái bắt tay tối đa trong cuộc họp là $C_n^2$ .\\
Trong thực tế thì tổng cộng số cái bắt tay là 2014; vì vậy $C_n^2>2014$\\
$\Rightarrow\dfrac{n!}{2\left(n-2\right)!}>2014\Rightarrow n\left(n-1\right)>4028\Rightarrow{n^2}-n-4028>0\Rightarrow n>63,97$ .
\begin{itemize}
\item Với $n=64$ thì số cái bắt tay tối đa là $C_{64}^2=2\,016$ ; số người mà ông Kim không bắt tay là $2016-2014=2$ (thỏa mãn).\\
\item Với $n=65$ thì số cái bắt tay tối đa là $C_{65}^2=2\,080$ ; số người mà ông Kim không bắt tay là $2080-2014=66$ (vô lí vì trong phòng họp đang có 65 người).\\
\item Ta không cần thử lại với $n>65$ vì luôn xảy ra điều vô lí như trên.
\end{itemize}
Vậy số người tham dự cuộc họp là $64$.}
\end{ex}
\begin{ex}%Câu 19
 \immini[thm]{Hình dáng phần đất liền của hai xã thuộc tỉnh Đồng Tháp được mô hình hóa bởi đồ thị hàm số $y=\dfrac{x^2+ax+b}{x-2}$ ; biết đồ thị có một điểm cực trị là $\left(1\,;\,\,1\right)$ , với hệ trục tọa độ Oxy như hình vẽ, đơn vị trên mỗi trục là 10 mét. Để thuận tiện cho giao thông hai xã, lãnh đạo tỉnh đã phê duyệt dự án xây một chiếc cầu nối phần đất liền của hai xã này. Nhằm tiết kiệm chi phí cho công trình, người kỹ sư trưởng thiết kế có nhiệm vụ nghiên cứu để chọn được hai vị trí A, B trên phần đất liền hai xã sao cho độ dài chiếc cầu (đoạn AB) là ngắn nhất có thể. Hỏi độ dài ngắn nhất của chiếc cầu đó (tính theo đường chim bay) là bao nhiêu mét (làm tròn đến hàng phần chục)?
 \shortans{43,9}}{\includegraphics[scale=1]{img/HXN-3.19}}
\loigiai{
Ta có $y'=\dfrac{x^2-4x-2a-b}{\left(x-2\right)^2}$ .\\
 Vì $\left(1\,;\,\,1\right)$ là điểm cực trị của đồ thị hàm số nên $\left\{\begin{aligned}
& y(1)=1\\ 
&{y}'(1)=0\\ 
\end{aligned}\right.\Rightarrow\left\{\begin{aligned}
& a+b=-2\\ 
& 2a+b=-3\\ 
\end{aligned}\right.\Rightarrow\left\{\begin{aligned}
& a=-1\\ 
& b=-1\\ 
\end{aligned}\right.$ .\\
Hàm số trở thành $y=\dfrac{x^2-x-1}{x-2}=x+1+\dfrac{1}{x-2},\,\,\,x\ne 2$ .\\
Gọi $A\left(2+a\,;\,\,3+a+\dfrac{1}{a}\right)\,,\,\,B\left(2-b\,;\,\,3-b-\dfrac{1}{b}\right)$ là hai điểm thuộc hai nhánh đồ thị với $a\,,\,\,b>0$ .\\
 Ta có: $A{B^2}=\left(a+b\right)^2+\left(a+b+\dfrac{1}{a}+\dfrac{1}{b}\right)^2=\left(a+b\right)^2+\left(a+b+\dfrac{a+b}{ab}\right)^2=\left(a+b\right)^2\left[1+\left(1+\dfrac{1}{ab}\right)^2\right]$\\
$=\left(a+b\right)^2\left(2+\dfrac{2}{ab}+\dfrac{1}{a^2b^2}\right)\overset{AM-GM}{\mathop{\ge}}\,4ab\left(2+\dfrac{2}{ab}+\dfrac{1}{a^2b^2}\right)$ $=8+8ab+\dfrac{4}{ab}\overset{AM-GM}{\mathop{\ge}}\,8+8\sqrt{2}$ .\\
Độ dài ngắn nhất của cây cầu (theo đường chim bay) là $AB\times 10=\sqrt{8+8\sqrt{2}}\times 10\approx\, 43,9 \,m$ .\\
Dấu “=” xảy ra khi và chỉ khi $a=b$ và $8ab=\dfrac{4}{ab}\Leftrightarrow a=b=\dfrac{1}{\sqrt[4]{2}}$ .}
\end{ex}
\begin{ex}%Câu 20
 \immini[thm]{Một cái chậu đựng nước có dạng hình chóp cụt đều đáy là các tam giác cạnh bằng 1 dm và 3 dm. Chiều cao chậu nước bằng 4 dm.
 Người ta bơm nước vào chậu với lưu lượng không đổi 0,5 lít/phút. Đến phút thứ 10 thì tốc độ dâng lên của nước trong chậu là bao nhiêu dm/phút? (Kết quả được làm tròn đến hàng phần trăm).
 \shortans{0,17}}{\includegraphics[scale=.8]{img/HXN-3.20}}
\loigiai{
\begin{center}
 \includegraphics[scale=1]{img/HXN-3.20a}
\end{center}
Gọi MN là độ dài cạnh tam giác đều theo mực nước tức thời (MN thay đổi).\\
Đặt $MN=a\times h+b$ (hàm số bậc nhất theo h).\\
Khi $h=0$ thì $MN=1$ ; khi $h=4$ thì $MN=3$ .\\
Do đó $\left\{\begin{aligned}
& b=1\\ 
& 4a+b=3\\ 
\end{aligned}\right.\Rightarrow\left\{\begin{aligned}
& b=1\\ 
& a=\dfrac{1}{2}\\ 
\end{aligned}\right.$ ; suy ra $MN=\dfrac{1}{2}h+1$ .\\
Diện tích mặt nước tức thời là $S=\dfrac{M{N^2}\sqrt{3}}{4}=\dfrac{\left(0,5h+1\right)^2\sqrt{3}}{4}$ .\\
Thể tích nước tức thời là $V=\dfrac{1}{3}h\left(S_0+\sqrt{S_0S}+S\right)$ ; $S_0$ là diện tích mặt nước ban đầu (đáy nhỏ).\\
$V=\dfrac{1}{3}h\left(\dfrac{\sqrt{3}}{4}+\sqrt{\dfrac{\sqrt{3}}{4}\cdot\dfrac{\left(0,5h+1\right)^2\sqrt{3}}{4}}+\dfrac{\left(0,5h+1\right)^2\sqrt{3}}{4}\right)$ $=\dfrac{h\sqrt{3}}{12}\left[1+\left(1+0,5h\right)+\left(1+0,5h\right)^2\right]$ .\\
$V=\dfrac{h\sqrt{3}}{12}\left(0,25h^2+1,5h+3\right)=\dfrac{\sqrt{3}}{12}\left(0,25h^3+1,5h^2+3h\right)$ (*).\\
Sau 10 phút thì lượng nước trong chậu là $V=0,5\times 10=5\,\,d{m^3}$ .\\
Do đó $\dfrac{\sqrt{3}}{12}\left(0,25h^3+1,5h^2+3h\right)=5\Rightarrow h\approx 3,27\,\,dm$ (Lưu vào A).\\
Từ (*) đạo hàm hai vế theo $t$ ta được:
$$ \dfrac{dV}{dt} = \dfrac{\sqrt{3}}{12} (0,75h^2 + 3h + 3) \cdot \dfrac{dh}{dt} \quad (**) $$
Thay các giá trị $h = A; V=5; \dfrac{dV}{dt}=0,5$ vào (**) ta được:
$$ \dfrac{dh}{dt} \approx 0,17 \ m/\text{phút} $$
Vậy tốc độ dâng lên của nước trong chậu xấp xỉ $0,17 \ m/\text{phút}$.
}
\end{ex}
\begin{ex}%Câu 21
Trong một trận đấu cờ vua của hai kỳ thủ là Lê Quang Liêm và vua cờ Carlsen; mỗi ván cờ luôn có kẻ thắng người thua (vì nếu hai kỳ thủ hòa thì sẽ bốc thăm để chọn người thắng ván đó). Biết rằng trong mỗi ván đấu, xác suất để anh Liêm dành chiến thắng bằng 0,4; xác suất để Carlsen dành chiến thắng bằng 0,6. Mỗi trận thắng được tính 1 điểm cho kỳ thủ, người thua không được điểm nào. Nếu người nào tạo được cách biệt 2 điểm thì sẽ dành chiến thắng chung cuộc. Tính xác suất để Lê Quang Liêm là người chiến thắng sau cùng (làm tròn đến hàng phần trăm).\shortans{0,31}
\loigiai{
Gọi $P(n)$ là xác suất để Lê Quang Liêm dành chiến thắng khi hiệu số điểm của anh so với Carlsen là $n$ điểm. Ta có $n \in \{-2; -1; 0; 1; 2\}$.
Ta cần tính xác suất chiến thắng của anh Liêm từ trạng thái $n=0$.
\\Theo công thức xác suất toàn phần: $P(0) = 0,4 \times P(1) + 0,6 \times P(-1)$ (1).
\\Ta có $P(2)=1$ và $P(-2)=0$; $P(1)=0,4 \times P(2) + 0,6 \times P(0)$ hay $P(1)=0,4+0,6 \times P(0)$ (2);
$P(-1)=0,4 \times P(0) + 0,6 \times P(-2)$ hay $P(-1)=0,4 \times P(0)$ (3).
\\Thay (2) và (3) vào (1): $P(0) = 0,4 [0,4 + 0,6 \times P(0)] + 0,6 \times 0,4 \times P(0) \Rightarrow P(0) = \dfrac{4}{13} \approx 0,31$.
\\Vậy xác suất để Lê Quang Liêm chiến thắng Carlsen là xấp xỉ $0,31$}
\end{ex}
\begin{ex}%Câu 22
Trong không gian Oxyz, cho ba điểm $A\left(0\,;\,3\,;\,-5\right),\,\,B\left(1\,;\,1\,;\,-5\right),\,\,C\left(4\,;\,3\,;\,-1\right)$ và mặt cầu $\left(S_m\right):{x^2}+y^2+z^2+\left(m-2\right)x+4y+\left(m-2\right)z-3=0$ (m là tham số thực). Gọi (T) là tập hợp tất cả điểm cố định mà mặt cầu $\left(S_m\right)$ luôn đi qua với mọi số thực m và M là một điểm di động trên (T) sao cho thể tích tứ diện MABC đạt giá trị lớn nhất $V_{\max}$ . Tính giá trị lớn nhất $V_{\max}$ đó (làm tròn đến hàng phần chục).
\shortans{15,3}
\loigiai{
 \begin{center}
 \includegraphics[scale=1]{img/HXN-3.22}
 \end{center}
Xét $M(x;y;z)$ là điểm mà $(S_m)$ luôn đi qua với mọi $m$.
\\Ta có: $x^2 + y^2 + z^2 + (m-2)x + 4y + (m-2)z - 3 = 0, \forall m \in \mathbb{R}$\\
$\Leftrightarrow m(x+z) + x^2 + y^2 + z^2 - 2x + 4y - 2z - 3 = 0, \forall m \in \mathbb{R}$\\
$\Leftrightarrow \begin{cases} x+z = 0 \\ x^2 + y^2 + z^2 - 2x + 4y - 2z - 3 = 0 \end{cases}$
\\Tập hợp điểm $M$ là đường tròn $(C)$ là giao tuyến của mặt phẳng $(P): x+z = 0$ và mặt cầu có tâm $I(1;-2;1)$, bán kính $R=3$;
\\Ta có: $d(I,(P)) = \dfrac{|1 \cdot 1 + 0 \cdot (-2) + 1 \cdot 1|}{\sqrt{1^2+0^2+1^2}} = \dfrac{|1+1|}{\sqrt{2}} = \dfrac{2}{\sqrt{2}} = \sqrt{2}$.
\\Vì vậy $(C)$ có bán kính là $r = \sqrt{R^2 - d^2(I,(P))} = \sqrt{7}$ và tâm $J(0;-2;0)$.
\\Mặt phẳng $(ABC)$ có phương trình $2x+y-2z-13=0$; ta thấy $(ABC)$ vuông góc với $(P)$.\\
$ \Rightarrow [\vec{AB},\vec{AC}] = (-8;-4;8) \Rightarrow S_{ABC} = \dfrac{1}{2}\sqrt{(-8)^2+(-4)^2+8^2} = 6$.
\\Thể tích tứ diện $MABC$ là $V_{MABC} = \dfrac{1}{3}d(M,(ABC)) \cdot S_{ABC} = \dfrac{1}{3}d(M,(ABC)) \cdot 6 = 2d(M,(ABC))$.
\\Thể tích này lớn nhất khi và chỉ khi $d(M,(ABC))$ đạt giá trị lớn nhất.
\\Ta có: $d(M,(ABC))_{\max} = r + d(J,(ABC)) = \sqrt{7} + \dfrac{|2(0) + 1(-2) - 2(0) - 13|}{\sqrt{2^2+1^2+(-2)^2}} = \sqrt{7} + 5$.
\\Vì vậy $V_{\max} = 2\sqrt{7} + 10 \approx 15,3$.}
\end{ex}
\Closesolutionfile{ans}
\inputansbox{6,4,3}{ans/ans-HXN-\sode-T,ans/ans-HXN-\sode-TF,ans/ans-HXN-\sode-SA}